% GENERATED FILE. Edit canonical lessons, then run npm run generate:books.
% canonical-source-hash: fnv1a64:e273aa137b285f45
% canonical-lessons: RU-C231-kopeyka, RU-C231-rubl, RU-C231-lektsiya, RU-C231-zanyatiye, RU-C231-otsenka, RU-R231-first-pass-words-from-chapters-227-228, RU-R231-second-pass-words-from-chapters-229-231

\chapter{Kopeck, Rouble, Lecture, Class, Mark}
\label{ch:ru-kopeck-rouble-lecture-class-mark}
\hlchaptermodality{231}

\begin{chapteropening}
\textbf{By the end of this chapter:} \emph{I can say a kopeck, a rouble, a lecture, a class and a mark.}

Complete the last lesson of chapter 231: I can say a kopeck, a rouble, a lecture, a class and a mark.
\end{chapteropening}

\section[\texorpdfstring{\ru{копейка} (kopéyka)}{копейка (kopéyka)}]{\ru{копейка} (kopéyka) — a kopeck}

\label{lesson:RU-C231-kopeyka}



\begin{quote}
Before the new one: say the Russian for a security guard, then the Russian for a nanny.
\end{quote}

\subsection*{You'll want to know: \ru{копейка}}
\textbf{\ru{копейка}} (\emph{kopéyka}) — \textquotedblleft{}a kopeck\textquotedblright{}.

\begin{grammarlens}[title={What you've built}]
One more word for this chapter.
\end{grammarlens}

\subsection*{Your turn}
\begin{itemize}
\raggedright
  \item \emph{Say it:} \emph{kopéyka}
  \item \emph{Say it:} \emph{kopéyka}, once more
  \item \emph{Say it:} say \emph{kopéyka}
  \item \emph{Recall:} say the Russian for a security guard, then the Russian for a nanny, then say \emph{kopéyka} again
\end{itemize}

\begin{tcolorbox}[breakable,colback=teal!4,colframe=teal!35!black,title={Before you move on}]
What does \ru{копейка} mean? (\textquotedblleft{}A kopeck\textquotedblright{}.) Say it once more.
\end{tcolorbox}
\section[\texorpdfstring{\ru{рубль} (rubl)}{рубль (rubl)}]{\ru{рубль} (rubl) — a rouble}

\label{lesson:RU-C231-rubl}



\begin{quote}
Before the new one: say the Russian for a nanny, then the Russian for a kopeck.
\end{quote}

\subsection*{You'll want to know: \ru{рубль}}
\textbf{\ru{рубль}} (\emph{rubl}) — \textquotedblleft{}a rouble\textquotedblright{}.

\begin{grammarlens}[title={What you've built}]
One more word for this chapter.
\end{grammarlens}

\subsection*{Your turn}
\begin{itemize}
\raggedright
  \item \emph{Say it:} \emph{rubl}
  \item \emph{Say it:} \emph{rubl}, once more
  \item \emph{Say it:} say \emph{rubl}
  \item \emph{Recall:} say the Russian for a nanny, then the Russian for a kopeck, then say \emph{rubl} again
\end{itemize}

\begin{tcolorbox}[breakable,colback=teal!4,colframe=teal!35!black,title={Before you move on}]
What does \ru{рубль} mean? (\textquotedblleft{}A rouble\textquotedblright{}.) Say it once more.
\end{tcolorbox}
\section[\texorpdfstring{\ru{лекция} (léktsiya)}{лекция (léktsiya)}]{\ru{лекция} (léktsiya) — a lecture}

\label{lesson:RU-C231-lektsiya}



\begin{quote}
Before the new one: say the Russian for a kopeck, then the Russian for a rouble.
\end{quote}

\subsection*{You'll want to know: \ru{лекция}}
\textbf{\ru{лекция}} (\emph{léktsiya}) — \textquotedblleft{}a lecture\textquotedblright{}.

\begin{grammarlens}[title={What you've built}]
One more word for this chapter.
\end{grammarlens}

\subsection*{Your turn}
\begin{itemize}
\raggedright
  \item \emph{Say it:} \emph{léktsiya}
  \item \emph{Say it:} \emph{léktsiya}, once more
  \item \emph{Say it:} say \emph{léktsiya}
  \item \emph{Recall:} say the Russian for a kopeck, then the Russian for a rouble, then say \emph{léktsiya} again
\end{itemize}

\begin{tcolorbox}[breakable,colback=teal!4,colframe=teal!35!black,title={Before you move on}]
What does \ru{лекция} mean? (\textquotedblleft{}A lecture\textquotedblright{}.) Say it once more.
\end{tcolorbox}
\section[\texorpdfstring{\ru{занятие} (zanyátiye)}{занятие (zanyátiye)}]{\ru{занятие} (zanyátiye) — a class, a lesson}

\label{lesson:RU-C231-zanyatiye}



\begin{quote}
Before the new one: say the Russian for a rouble, then the Russian for a lecture.
\end{quote}

\subsection*{You'll want to know: \ru{занятие}}
\textbf{\ru{занятие}} (\emph{zanyátiye}) — \textquotedblleft{}a class, a lesson\textquotedblright{}.

\begin{grammarlens}[title={What you've built}]
One more word for this chapter.
\end{grammarlens}

\subsection*{Your turn}
\begin{itemize}
\raggedright
  \item \emph{Say it:} \emph{zanyátiye}
  \item \emph{Say it:} \emph{zanyátiye}, once more
  \item \emph{Say it:} say \emph{zanyátiye}
  \item \emph{Recall:} say the Russian for a rouble, then the Russian for a lecture, then say \emph{zanyátiye} again
\end{itemize}

\begin{tcolorbox}[breakable,colback=teal!4,colframe=teal!35!black,title={Before you move on}]
What does \ru{занятие} mean? (\textquotedblleft{}A class, a lesson\textquotedblright{}.) Say it once more.
\end{tcolorbox}
\section[\texorpdfstring{\ru{оценка} (otsénka)}{оценка (otsénka)}]{\ru{оценка} (otsénka) — a mark, a grade}

\label{lesson:RU-C231-otsenka}



\begin{quote}
Before the new one: say the Russian for a lecture, then the Russian for a class.
\end{quote}

\subsection*{You'll want to know: \ru{оценка}}
\textbf{\ru{оценка}} (\emph{otsénka}) — \textquotedblleft{}a mark, a grade\textquotedblright{}.

\begin{grammarlens}[title={What you've built}]
One more word for this chapter.
\end{grammarlens}

\subsection*{Your turn}
\begin{itemize}
\raggedright
  \item \emph{Say it:} \emph{otsénka}
  \item \emph{Say it:} \emph{otsénka}, once more
  \item \emph{Say it:} say \emph{otsénka}
  \item \emph{Recall:} say the Russian for a lecture, then the Russian for a class, then say \emph{otsénka} again
\end{itemize}

\begin{tcolorbox}[breakable,colback=teal!4,colframe=teal!35!black,title={Before you move on}]
What does \ru{оценка} mean? (\textquotedblleft{}A mark, a grade\textquotedblright{}.) Say it once more.
\end{tcolorbox}
\section[(dialogue)]{Words from chapters 227-228 — a first pass}

\label{lesson:RU-R231-first-pass-words-from-chapters-227-228}



\begin{quote}
Say the words for a class and a mark.
\end{quote}

\subsection*{Your turn}
Say these aloud:

\begin{itemize}
\raggedright
  \item a fine, a person, a man, a woman, a boy
  \item a mistake, a rule, an example, a meaning, a first name
\end{itemize}

\begin{tcolorbox}[breakable,colback=teal!4,colframe=teal!35!black,title={Before you move on}]
A fine? (\textbf{\ru{штраф}}.) A person? (\textbf{\ru{человек}}.) A man? (\textbf{\ru{мужчина}}.) A woman? (\textbf{\ru{женщина}}.) A boy? (\textbf{\ru{мальчик}}.) A mistake? (\textbf{\ru{ошибка}}.) A rule? (\textbf{\ru{правило}}.) An example? (\textbf{\ru{пример}}.) A meaning? (\textbf{\ru{значение}}.) A first name? (\textbf{\ru{имя}}.)
\end{tcolorbox}
\section[(dialogue)]{Words from chapters 229-231 — a second pass}

\label{lesson:RU-R231-second-pass-words-from-chapters-229-231}



\begin{quote}
Say the words for a secretary and an accountant.
\end{quote}

\subsection*{Your turn}
Say these aloud:

\begin{itemize}
\raggedright
  \item a secretary, an accountant, a manager, a programmer, a lecturer
  \item a gardener, a cashier, a hairdresser, a security guard, a nanny
  \item a kopeck, a rouble, a lecture, a class, a mark
\end{itemize}

\begin{tcolorbox}[breakable,colback=teal!4,colframe=teal!35!black,title={Before you move on}]
A secretary? (\textbf{\ru{секретарь}}.) An accountant? (\textbf{\ru{бухгалтер}}.) A manager? (\textbf{\ru{менеджер}}.) A programmer? (\textbf{\ru{программист}}.) A lecturer? (\textbf{\ru{преподаватель}}.) A gardener? (\textbf{\ru{садовник}}.) A cashier? (\textbf{\ru{кассир}}.) A hairdresser? (\textbf{\ru{парикмахер}}.) A security guard? (\textbf{\ru{охранник}}.) A nanny? (\textbf{\ru{няня}}.) A kopeck? (\textbf{\ru{копейка}}.) A rouble? (\textbf{\ru{рубль}}.) A lecture? (\textbf{\ru{лекция}}.) A class? (\textbf{\ru{занятие}}.) A mark? (\textbf{\ru{оценка}}.)
\end{tcolorbox}
